% New proof fragment. No canonical source has been edited by this lane.
% Uses the actual period construction, finite normalization certificate,
% toric fan, and analytically proved canonical-pencil recovery.
\subsection{Exact global identifications of the retained period parameter}
\label{pshift:sec:classification}

Keep the original base $B$, coordinate $t$, values
$t(p_1)=0,t(p_2)=1,t(p_0)=\infty$, and the original functions
$\tau,\mu,\beta^{\mathrm{part}}$. Write
\[
 \beta_c=\beta^{\mathrm{part}}+c,\qquad
 \mathcal P=\{c\in\mathbb C:\operatorname{Im}c<-M\},\qquad
 \Pi_c(z)=\begin{pmatrix}6\mu(z)&\tau(z)&1&0\\
                      \beta_c(z)&\mu(z)&0&1\end{pmatrix}.
 \tag{PS1}\label{pshift:eq:periods}
\]
All assertions below concern the analytic pieces and their actual gluing.
No integral cohomology, sphere recognition, or CDP assertion is an input.
In particular, the global parameter classification is not inferred from
an equality between two fibre lattices.

Put $e=(0,1)^{\mathsf t}\in\mathbb C^2$ and retain the ordered lattice
\[
 \Lambda=\mathbb Z\langle\widehat\gamma,\widehat u,
                        \widehat w,\widehat\delta\rangle,
 \qquad E=\widehat\delta\otimes\gamma.
 \tag{PS2}\label{pshift:eq:rank-one-map}
\]
Thus $E$ is the matrix with its only nonzero entry, equal to one, in
row four and column one; $E^2=0$ and $\Pi_cE=e\gamma$.
The finite data, with every sign retained, are
\[
\begin{gathered}
 A_1=\begin{pmatrix}1&0&0&0\\6&0&1&0\\-6&-1&-1&0\\-2&1&0&1\end{pmatrix},
 \qquad
 A_2=\begin{pmatrix}1&0&0&0\\0&0&-1&0\\-6&1&0&0\\3&0&1&1\end{pmatrix},\\
 v_1=(1,2,-4,0)^{\mathsf t},\quad
 v_2=(-1,-3,3,0)^{\mathsf t},\quad
 (m_1,m_2)=(3,4),\quad
 (\epsilon_1,\epsilon_2)=(1,-1),\quad \epsilon_j=\gamma(v_j).
\end{gathered}
\tag{PS3}\label{pshift:eq:finite-data}
\]
The covariance is $R_g\Pi_c=\Pi_c(gz)A_g$, with
\[
 R_{g_1}=\begin{pmatrix}-1/\tau&0\\(1-\mu)/\tau&1\end{pmatrix},\qquad
 R_{g_2}=\begin{pmatrix}1/\tau&0\\-\mu/\tau&1\end{pmatrix},\qquad
 R_{g_0}=I_2.
 \tag{PS4}\label{pshift:eq:covariance}
\]
In particular $R_ge=e$ for every $g$.

\begin{theorem}[The exact identification subgroup is $2\mathbb Z$]
\label{pshift:thm:classification}
For $c,c'\in\mathcal P$, the retained analytic threefolds are
biholomorphic if and only if
\[
                         c'-c\in2\mathbb Z.
 \tag{PS5}\label{pshift:eq:classification}
\]
Every such biholomorphism preserves the original map to $B$. For every
even integer $n$, the proof below constructs an isomorphism
$X_c\longrightarrow X_{c+n}$ whose complex-linear part on the original
universal fibre coordinates is $I_2$. Its map on the marked fibre lattice
is $I-nE$.

The equality of the middle lattices for all integer shifts loses a
nonzero order-two finite-filling invariant: the shift $n$ changes that
invariant by $n\bmod2$. The cusp contributes no additional obstruction
to the displayed even-shift construction.
\end{theorem}

\subsubsection{The middle family and its complete rational centralizer}

For an integer $n$, put $S_n=I+nE$. Since $S_n^{-1}=S_{-n}$,
\[
 \Pi_{c+n}=\Pi_c S_n,\qquad
 \Pi_{c+n}\Lambda=\Pi_c\Lambda.
 \tag{PS6}\label{pshift:eq:integer-shear}
\]
Both original matrices fix $\widehat\delta$ and preserve $\gamma$.
Consequently $A_jE=EA_j=E$, and the same commutation holds for every
monodromy word. The identity of $\mathbb C^2$ therefore identifies
the two full middle torus families with their actual $R_g$-action.
Its source-to-target lattice map is $S_{-n}$, since
$\Pi_{c+n}S_{-n}=\Pi_c$.

\begin{lemma}[All middle-family linear parts]
\label{pshift:lem:centralizer}
The simultaneous rational centralizer of $A_1,A_2$ is exactly
\[
 \{aI+bE:a,b\in\mathbb Q\}.
 \tag{PS7}\label{pshift:eq:centralizer}
\]
A biholomorphism between the middle families for $c,c'$ over the
identity of $B^\circ$ forces $c'-c\in\mathbb Z$. Its complex-linear
part on every fibre is either $I_2$ or $-I_2$.
\end{lemma}

\begin{proof}
Solving the two fixed-vector systems from (PS3) gives
\[
 (\Lambda\otimes\mathbb Q)^{A_1}
   =\mathbb Q(1,2,-4,0)^{\mathsf t}+\mathbb Q\widehat\delta,
 \quad
 (\Lambda\otimes\mathbb Q)^{A_2}
   =\mathbb Q(1,3,-3,0)^{\mathsf t}+\mathbb Q\widehat\delta.
\]
Their intersection is $\mathbb Q\widehat\delta$. The fixed covectors
are respectively
\[
 \mathbb Q\gamma+\mathbb Q(2u+w+3\delta),\qquad
 \mathbb Q\gamma+\mathbb Q(u+w+2\delta).
\]
Their intersection is $\mathbb Q\gamma$: comparing the $u$ and $w$
coefficients of a common multiple gives $2x=y$ and $x=y$, hence
$x=y=0$. Every commuting endomorphism therefore preserves
$\mathbb Q\widehat\delta$ and $\ker\gamma$.

On $\ker\gamma/\mathbb Q\widehat\delta$, in the classes of
$(\widehat u,\widehat w)$, the two matrices are
\[
 Q_1=\begin{pmatrix}0&1\\-1&-1\end{pmatrix},\qquad
 Q_2=\begin{pmatrix}0&-1\\1&0\end{pmatrix}.
\]
A matrix commuting with $Q_2$ has the form
$\left(\begin{smallmatrix}b&d\\-d&b\end{smallmatrix}\right)$.
Commutation with $Q_1$ gives $d=0$. Thus a commuting endomorphism $H$
has columns
\[
 H\widehat\gamma=(a,r,s,k)^{\mathsf t},\quad
 H\widehat u=b\widehat u+u_0\widehat\delta,\quad
 H\widehat w=b\widehat w+w_0\widehat\delta,\quad
 H\widehat\delta=d_0\widehat\delta.
\]
Use the actual columns $A_1\widehat u=-\widehat w+\widehat\delta$,
$A_1\widehat w=\widehat u-\widehat w$, and
$A_2\widehat u=\widehat w$. Commutation on these columns gives
\[
 d_0-w_0=b+u_0,\qquad u_0=2w_0,\qquad w_0=u_0.
\]
Hence $u_0=w_0=0$ and $d_0=b$. Commutation on
$\widehat\gamma$ with $A_1$ gives
\[
 r=2(a-b),\qquad s=-4(a-b).
\]
Commutation on that column with $A_2$ also gives $r=-s$.
Therefore $a=b$, $r=s=0$, and $H=aI+kE$. Conversely these matrices
commute, by the fixed-vector and fixed-covector identities preceding
(PS6). This proves (PS7).

For the geometric assertion, a holomorphic map between two compact
complex tori lifts to their universal covers. The derivative of the
lift is periodic, so its entries descend to holomorphic functions on
the compact source torus and are constant. The lift is consequently
affine, with a complex-linear part $L$. In a holomorphic family these
linear parts vary holomorphically. Their lattice maps are integral
and locally constant, hence give an isomorphism of the two marked
local systems. Because the map is over the identity of the base,
that integral map $H$ commutes with both $A_j$. Thus
$H=aI+bE$, with $a,b\in\mathbb Z$. Its determinant is $a^4$;
invertibility over $\mathbb Z$ gives $a=1$ or $a=-1$.

The complete period equation is
\[
                   L(z)\Pi_c(z)=\Pi_{c'}(z)(aI+bE).
 \tag{PS8}\label{pshift:eq:period-morphism}
\]
Its third and fourth columns give $L(z)=aI_2$. Its first column
then gives $a(c'-c)+b=0$. Therefore $c'-c=-b/a\in\mathbb Z$,
and the asserted two linear parts exhaust the possibilities.
\end{proof}

\subsubsection{The exact finite obstruction, including normalization}

The two cyclic norm matrices are, by multiplication of (PS3),
\[
 N_1=I+A_1+A_1^2=
 \begin{pmatrix}3&0&0&0\\6&0&0&0\\-12&0&0&0\\0&2&1&3\end{pmatrix},
\]
\[
 N_2=I+A_2+A_2^2+A_2^3=
 \begin{pmatrix}4&0&0&0\\12&0&0&0\\-12&0&0&0\\0&2&2&4\end{pmatrix}.
 \tag{PS9}\label{pshift:eq:norms}
\]
For example these formulas give
$N_1\widehat w=\widehat\delta$,
$N_2\widehat u=N_2\widehat w=2\widehat\delta$,
$N_1\widehat\delta=3\widehat\delta$, and
$N_2\widehat\delta=4\widehat\delta$.

We first justify why a local isomorphism must satisfy the exact affine
equations on the original finite cover. The actual ramified base change
$t_j=u^{m_j}$ has local equation
\[
 s_j^{m_j}-u^{m_j}
   =\prod_{a=0}^{m_j-1}(s_j-e^{2\pi ia/m_j}u)=0.
\]
Its normalization separates all these smooth branches, each with
coordinates $(z_1,z_2,u)$. The original map
$y\mapsto(q_j(y),s_j(y))$ identifies the original torus family $J'_j$
with this normalization, as proved by the explicit branches in
\cref{ffcert:thm:normalization}. A biholomorphism over the identity of
$D_j$ base-changes over the identity of the $u$-disc and lifts uniquely
to the normalization: on local rings it carries integral elements in
the total fraction ring to integral elements, and its inverse does
the same. This lift commutes with the original deck action because
both composites lift the same map $(x,u)\mapsto(F(x),\zeta_j u)$.
In particular one may not replace the generator by an arbitrary
coprime power.

Let $c'=c+n$, where $n$ is an integer, and suppose the linear part of
the lifted map is $aI_2$, $a\in\{1,-1\}$. The lifted map on the torus
family is $[z,\zeta]\mapsto[z,a\zeta+B_j(z)]$. Its translating
section has a holomorphic lift $B_j:\Delta_j\to\mathbb C^2$:
$\Delta_j\times\mathbb C^2\to J'_j$ is a covering map and the disc is
simply connected. Local inverses of this covering show that the lift
is holomorphic. The deck-equivariance identity holds modulo an
actual lattice vector, so, after choosing this lift, it is
\[
 B_j(g_jz)-R_{g_j}(z)B_j(z)
  =\Pi_c(g_jz)\left(
       \frac{(1-a)v_j+n\epsilon_j\widehat\delta}{m_j}
          +\ell_j\right),\qquad \ell_j\in\Lambda.
 \tag{PS10}\label{pshift:eq:finite-equation}
\]
The integer column $\ell_j$ is constant: it is locally constant by
discreteness of the marked lattice and the disc is connected.
Iterating $m_j$ times, with $A_jv_j=v_j$ and
$A_j\widehat\delta=\widehat\delta$, yields the necessary equality
\[
 N_j\ell_j=-(1-a)v_j-n\epsilon_j\widehat\delta.
 \tag{PS11}\label{pshift:eq:finite-norm-obstruction}
\]
Here the left side of the iterated equation is zero because the
base point and the original vector coordinate return after $m_j$
steps; real invertibility of $\Pi_c$ makes the lattice equality exact.

If $a=-1$, applying $\gamma$ to (PS11) gives
\[
                   m_j\gamma(\ell_j)=-2\epsilon_j.
\]
This is impossible already at $j=1$, since $3$ does not divide $2$.
Thus a global isomorphism cannot have negative linear part.

For $a=1$, the first row of (PS11) gives $\gamma(\ell_j)=0$.
Equation (PS9) then proves that the order-three filling imposes no
restriction on $n$, whereas the order-four filling imposes precisely
$n\in2\mathbb Z$. More intrinsically, the finite obstruction is the
class of $-n\epsilon_j\widehat\delta$ in
\[
 \Lambda^{A_j}/N_j\Lambda.
 \tag{PS12}\label{pshift:eq:finite-class}
\]
In the original fixed bases this quotient is
$\mathbb Z/3$ at $j=1$, with $\widehat\delta$ zero, and
$\mathbb Z/4\oplus\mathbb Z/2$ at $j=2$, with
$\widehat\delta$ generating the second summand. Indeed (PS9) gives
\[
 N_1\Lambda=3\mathbb Z v_1\oplus\mathbb Z\widehat\delta,
 \quad
 N_2\Lambda=4\mathbb Z\varepsilon'\oplus
                      2\mathbb Z\widehat\delta,
 \qquad \varepsilon'=(1,3,-3,0)^{\mathsf t}.
\]
This proves the exact parity obstruction rather than merely failure
of the untranslated identity.

\subsubsection{Local maps for every even shift}

Let $n\in2\mathbb Z$. Define the following original rational columns,
integral cocycles, and holomorphic translation vectors:
\[
\begin{aligned}
 r_1&=-\frac n3(\widehat u+\widehat w),&
 \ell_1&=-n\widehat w,&
 b_1(z)&=\Pi_c(z)r_1=-\frac n3(\tau(z)+1,\mu(z))^{\mathsf t},\\
 r_2&= \frac n4(\widehat u+\widehat w),&
 \ell_2&=\frac n2\widehat u,&
 b_2(z)&=\Pi_c(z)r_2=\frac n4(\tau(z)+1,\mu(z))^{\mathsf t}.
\end{aligned}
\tag{PS13}\label{pshift:eq:local-corrections}
\]
Direct multiplication of the complete matrices in (PS3) gives
\[
 (I-A_j)r_j=\frac{n\epsilon_j}{m_j}\widehat\delta+\ell_j.
 \tag{PS14}\label{pshift:eq:local-matrix-identity}
\]
For $j=1$ the four entries are $(0,0,-n,n/3)$, and for $j=2$
they are $(0,n/2,0,-n/4)$. Pairing with the full $\Pi_c(g_jz)$
shows that $b_j$ satisfies (PS10) for $a=1$.
Translation by $b_j$ is therefore already an isomorphism of the
two finite local fillings. The remaining issue is its exact overlap
with the fixed middle identification.

Keep the original linearizing coordinate, radius, and phase
$s_j(g_jz)=e^{-2\pi i/m_j}s_j(z)$. On a logarithm branch put
\[
 w_j(z)=\frac{\log s_j(z)}{2\pi i},\qquad
 a_j(z)=b_j(z)+n\epsilon_j w_j(z)e.
 \tag{PS15}\label{pshift:eq:finite-middle-sections}
\]
Changing the logarithm adds $n\epsilon_j\widehat\delta$ under the
period map, so $a_j$ is branch independent as a torus section.
For the specified negative generator, $w_j(g_jz)=w_j(z)-1/m_j$.
Equations (PS4) and (PS14) consequently give the exact equality
\[
 a_j(g_jz)=R_{g_j}(z)a_j(z)+\Pi_c(g_jz)\ell_j.
 \tag{PS16}\label{pshift:eq:local-affine-cocycle}
\]
This is why the logarithmic term in (PS15) cannot be omitted.
In fact the original finite overlap maps are translations by
$\sigma_{j,c}=w_j\Pi_c v_j$, and
\[
                  \sigma_{j,c+n}-\sigma_{j,c}
                     =n\epsilon_jw_je.
 \tag{PS17}\label{pshift:eq:actual-overlap-difference}
\]
Thus $a_j$ is exactly the middle translation induced by the local
translation $b_j$, with the original gluing maps retained.

At the cusp retain $s=\tau-h$, $e^{2\pi is}=t_c$, $s(g_0z)=s(z)-1$,
and the complete period matrix
\[
 Z_c=sB_0+C_c(t_c),\quad
 B_0=\begin{pmatrix}0&1\\-1&0\end{pmatrix},\quad
 C_c=\begin{pmatrix}6\mu&h\\ b^{\mathrm{part}}-h+c&\mu\end{pmatrix}.
 \tag{PS18}\label{pshift:eq:cusp-periods}
\]
For $\lambda=(\lambda_1,\lambda_2)^{\mathsf t}\in\mathbb Z^2$,
$C_{c+n}\lambda-C_c\lambda=(0,n\lambda_1)^{\mathsf t}$. Therefore
\[
 \exp(2\pi iC_{c+n}\lambda)t_c^{B_0\lambda}
       =\exp(2\pi iC_c\lambda)t_c^{B_0\lambda}.
 \tag{PS19}\label{pshift:eq:cusp-actions-identical}
\]
The entire cusp actions, including their nonconstant holomorphic units,
are identical for an integer shift.

The affine cocycle convention in (PS16) has multiplication law
\[
                    \ell_{gh}=\ell_g+A_g\ell_h.
 \tag{PS20}\label{pshift:eq:cocycle-law}
\]
Since $g_0=(g_1g_2)^{-1}$, it gives
\[
\begin{aligned}
 \ell_1+A_1\ell_2
   &=-\frac{3n}{2}\widehat w+\frac n2\widehat\delta,\\
 \ell_0=-A_{g_0}(\ell_1+A_1\ell_2)
   &=\frac{3n}{2}\widehat w-\frac n2\widehat\delta.
\end{aligned}
\tag{PS21}\label{pshift:eq:cusp-cocycle}
\]
Here $A_{g_0}$ fixes both toric lattice vectors, as is also immediate
from $A_{g_0}=\left(\begin{smallmatrix}I&0\\B_0&I\end{smallmatrix}\right)$.
Put
\[
 q_n=(-3n/2,n/2)^{\mathsf t}\in\mathbb Z^2,
                  \qquad a_0(z)=s(z)q_n.
 \tag{PS22}\label{pshift:eq:cusp-translation}
\]
Then $a_0(g_0z)-a_0(z)=-q_n=\Pi_c(g_0z)\ell_0$, exactly as
required by (PS16). Exponentiating gives the map
$(x,t_c)\mapsto(t_c^{q_n}x,t_c)$ on the dense cusp torus.
It extends holomorphically and invertibly to the original infinite fan:
the integral map
\[
 (y,y_3)\longmapsto(y+y_3q_n,y_3)
\]
translates each height-one triangle by the integer vector $q_n$ and
therefore permutes the actual cones, with inverse given by $-q_n$.
It commutes with the cusp quotient action in (PS19). This proves
extension across every cusp stratum, not only across its dense torus.

\subsubsection{The global additive correction and all overlap identities}

We give the global correction explicitly as an additive gluing
construction and prove its solvability. Let $\widetilde B^\circ$ be
the universal cover of the three-punctured base. Its fundamental group
is free on the original two negative peripheral generators. Pull back
the original $\Pi_c,R$ to this cover. Assign $\ell_1,\ell_2$ from
(PS13) to its free generators and extend by (PS20). Freeness makes
this an actual integral cocycle; its third peripheral value is
precisely (PS21). The maps
\[
 (z,y)\longmapsto
       (gz,R_g(z)y+\Pi_c(gz)\ell_g)
 \tag{PS23}\label{pshift:eq:affine-bundle-action}
\]
form an affine action. To check the action law, covariance gives
$R_g(hz)\Pi_c(hz)=\Pi_c(ghz)A_g$; substituting this equality
in the composition gives exactly (PS20).
The quotient is an affine holomorphic bundle over $B^\circ$, since
the deck action on the base is free and properly discontinuous and
one may use evenly covered base opens as affine charts.

On a finite punctured disc, (PS16) says that $a_j$ is a section of
this affine bundle. On the cusp punctured disc, (PS21)--(PS22) say
the same for $a_0$. Define a sheaf $\mathscr H_n$ on $B$ as follows.
Away from the special points it is the sheaf of sections of (PS23).
Near $p_j$, a section is required to have the form $a_j+\eta_j$,
where $\eta_j:\Delta_j\to\mathbb C^2$ is holomorphic and obeys
$\eta_j(g_jz)=R_{g_j}(z)\eta_j(z)$. Near the cusp it is required
to have the form $a_0+\eta_0$, with $\eta_0(t_c)$ holomorphic on
the full cusp disc. These are compatible definitions on the punctured
overlaps by the equal affine cocycles just proved. They also prove
local nonemptiness, including at all three special points.

The differences of these local sections form a vector sheaf
$\mathscr V$. Its ordinary sections have covariance $\eta(gz)=R_g\eta(z)$;
at the two finite points and the cusp its extensions are precisely
the holomorphic extensions specified above. Projection to the first
coordinate gives an exact sequence on the original compact base
\[
 0\longrightarrow\mathcal O_B e\longrightarrow\mathscr V
        \longrightarrow\mathcal L\longrightarrow0,
                  \qquad \mathcal L\simeq\mathcal O_B(-p_0).
 \tag{PS24}\label{pshift:eq:vector-sheaf}
\]
Here $\mathcal L$ is exactly the homogeneous period sheaf from
\cref{prop:period-input-audit}, not a newly chosen extension. Indeed
the first diagonal entries in (PS4) are its factors $-1/\tau$ and
$1/\tau$, and at the cusp both extensions require ordinary
holomorphicity. The original function
\[
 \Theta=\frac{(E_4\circ\tau)^2(E_6\circ\tau)^{1/2}}
                 {\Delta_{\mathrm{mod}}\circ\tau}
\]
has those exact factors, orders two and one in $s_1,s_2$, and a
simple pole in $t_c$. Dividing by $\Theta$ identifies the ordinary
and finite stalks with $\mathcal O_B$, and the cusp stalk with
$\mathfrak m_{p_0}$. This recovers the displayed $\mathcal L$
with its original cusp condition. These transformation and order
calculations are the ones used to construct the retained $\mu$;
no topological conclusion about $X_c$ is involved.

For completeness, exactness on the finite stalks in (PS24) retains
the off-diagonal entry of $R_g$. Given an equivariant first component
$u(z)$, begin with the vector $(u(z),0)^{\mathsf t}$ and average it
under the finite group action on vector-valued holomorphic functions.
The resulting vector is equivariant, and its first component is
still $u$. Thus projection is surjective there. Its kernel is
exactly the invariant second-component functions, which are
holomorphic functions of $s_j^{m_j}=t_j$, multiplied by $e$.
The ordinary and cusp stalks are immediate. This also proves that
$\mathscr V$ is locally free of rank two: lift a local frame of
$\mathcal L$ and join it to $e$.

Both $H^1(B,\mathcal O_B)$ and
$H^1(B,\mathcal O_B(-p_0))$ vanish. One can compute these two
groups without a degree approximation: on the two standard affine
charts of $\mathbb P^1$, split a holomorphic Laurent series on
their overlap into its nonnegative and negative powers. For
$\mathcal O_B$ these are the two coboundary contributions, and
for $\mathcal O_B(-p_0)$ its transition factor shifts the second
contribution by one, still leaving no missing exponent. The positive
part converges on the full first affine chart and the transformed
negative part on the full second affine chart. This proves the two
vanishings, equivalently the usual additive Cousin calculation on
these charts. The exact sequence (PS24) now gives
\[
                             H^1(B,\mathscr V)=0.
 \tag{PS25}\label{pshift:eq:no-analytic-obstruction}
\]

Here is how to choose the correction entering the maps. Choose
local sections $h_\alpha$ of $\mathscr H_n$, taking the displayed
$a_j,a_0$ on sufficiently small special discs and affine charts
elsewhere. Their differences $d_{\alpha\beta}=h_\beta-h_\alpha$
are a Cech $1$-cocycle with values in $\mathscr V$. Project it to
$\mathcal L$ and solve that additive cocycle, using its just proved
$H^1$ vanishing. Lift the resulting local $\mathcal L$ sections to
$\mathscr V$ by the local lifting construction in (PS24). Subtracting
their coboundary from $d$ leaves a scalar $\mathcal O_Be$ cocycle.
Solve that scalar additive cocycle by its $H^1$ vanishing. This
produces sections $v_\alpha$ with
$d_{\alpha\beta}=v_\beta-v_\alpha$. The local sections
$h_\alpha-v_\alpha$ agree and define a global section
$a\in\mathscr H_n(B)$. This is a two-stage construction of the
required correction, not an unproved additional gluing condition.
It gives holomorphic vectors $\eta_j,\eta_0$ with
\[
 a=a_j+\eta_j\quad(j=1,2),\qquad a=a_0+\eta_0
                              \quad\hbox{at the cusp}.
 \tag{PS26}\label{pshift:eq:global-correction}
\]
Any choices made in the additive solutions give maps with the same
asserted properties. Their differences as sections of this affine
bundle lie in $H^0(B,\mathscr V)=\mathbb C e$: this follows from
(PS24), $H^0(B,\mathcal L)=0$, and $H^0(B,\mathcal O_B)=\mathbb C$.

Projecting $a$ modulo $\Pi_c\Lambda$ gives a holomorphic section of
the middle torus family because every $\ell_g$ in (PS23) is integral.
The promised three kinds of local maps are therefore
\[
\begin{aligned}
 F_{\mathrm{mid}}[z,\zeta]&=[z,\zeta+a(z)],\\
 F_j[z,\zeta]&=[z,\zeta+b_j(z)+\eta_j(z)]\quad(j=1,2),\\
 F_0[x,t_c]&=[\exp(2\pi i\eta_0(t_c))t_c^{q_n}x,t_c].
\end{aligned}
\tag{PS27}\label{pshift:eq:full-isomorphism}
\]
The first descends by (PS23) and equality of the integer-shift
lattices. The finite maps descend by (PS10), (PS14), and the exact
covariance of $\eta_j$. The cusp map descends by (PS19); its
holomorphic-unit factor extends by torus multiplication on every
original chart, and its integer factor extends by the fan map in
(PS22). Each map is a biholomorphism, with inverse given by the
negative translation or the inverse torus multiplier and fan map.

On a finite overlap, write the original gluing as
$h_{j,c}(\zeta)=\zeta+\sigma_{j,c}$. Then (PS17), (PS26), and
(PS27) give the equality of actual torus maps
\[
                     h_{j,c+n}F_j=F_{\mathrm{mid}}h_{j,c}.
 \tag{PS28}\label{pshift:eq:finite-overlap-identity}
\]
Both sides add
$\sigma_{j,c}+a=\sigma_{j,c+n}+b_j+\eta_j$.
On the cusp overlap the original map is exponentiation of $\zeta$;
(PS22) and (PS26) give
$\exp(2\pi ia)=t_c^{q_n}\exp(2\pi i\eta_0)$, exactly the
third line of (PS27). Thus that overlap also agrees, including
all nonconstant period units. The maps and their inverses glue to
a global biholomorphism $X_c\to X_{c+n}$ over $B$.

The choices can be made coherently for every parameter and every even
shift. Both initial cocycle columns in (PS13) have $\gamma=0$.
Because every $A_g$ preserves $\gamma$, (PS20) shows that
$\gamma(\ell_g)=0$ for every word. Therefore the affine term
$\Pi_c(gz)\ell_g$ in (PS23) is independent of $c$. The displayed
$a_j,a_0$, and the vector sheaf $\mathscr V$ are also independent
of $c$. Perform the additive construction once for $n=2$, keeping
one resulting section $a^{(2)}$. For $n=2k$ choose
$a^{(2k)}=k a^{(2)}$ and multiply its local correcting vectors by
$k$. All their affine cocycles and local formulas are then exactly
those for $n=2k$. The resulting maps satisfy the full group law
\[
 F_{2h,c+2k}\circ F_{2k,c}=F_{2(h+k),c},\qquad
                 F_{0,c}=\mathrm{id}_{X_c}.
 \tag{PS28a}\label{pshift:eq:coherent-shift-action}
\]
This follows by adding the actual translation vectors in the first
two lines of (PS27), and by multiplying their units and adding the
integer exponents in its third line. The formulas are jointly
holomorphic in $c$, since their corrections are independent of it.

\subsubsection{Intrinsic necessity and the resulting parameter space}

The analytically computed complete linear system
$|\omega_{X_c}^{-2}|$ is the original map $f_c$, by
\cref{thm:canonical-ring-recovery}. A biholomorphism preserves its
canonical bundle, the complete space of sections, and its evaluation
map. Consequently it induces $\alpha\in\operatorname{PGL}_2(\mathbb C)$
with $f_{c'}F=\alpha f_c$. The three special fibres have distinct
intrinsic types: the cusp is reduced and nonnormal, while the other
two have multiplicities three and four. All remaining fibres are
smooth reduced tori. Thus $\alpha$ fixes $p_0,p_1,p_2$ individually.
In the original coordinate, a fractional-linear map fixing
$\infty,0,1$ has zero lower-left and upper-right entries and equal
diagonal entries; hence $\alpha=\mathrm{id}_B$.

Now \cref{pshift:lem:centralizer} forces $c'-c=n\in\mathbb Z$
and linear part $\pm I_2$. Equations (PS10)--(PS11), obtained on
the original normalized finite covers, exclude $-I_2$ and then
exclude odd $n$ for $+I_2$. Conversely (PS27)--(PS28) construct
the global map for every even $n$. This proves
\cref{pshift:thm:classification} without converting a failed
particular map into a nonisomorphism assertion.

As a complex parameter space for the isomorphism classes represented
by this retained family, the quotient $\mathcal P/(2\mathbb Z)$ has
the explicit coordinate
\[
 w=\exp(-\pi ic),\qquad 0<|w|<\exp(-\pi M).
 \tag{PS29}\label{pshift:eq:parameter-disc}
\]
Indeed $|w|=\exp(\pi\operatorname{Im}c)$, and two complex numbers
have the same value of this exponential exactly when they differ
by an even integer. Every point in the displayed punctured disc
has a logarithm giving a $c\in\mathcal P$. The exponential has
nonzero derivative, so these local inverse logarithms give the
claimed quotient complex structure. The integer-period quotient
seen by the middle family instead has coordinate $w^2$; it is
covered two-to-one by (PS29). Its two lifts are distinguished by
the exact order-two summand in (PS12). This statement classifies
the parameters of the retained family; it does not assert that
the family exhausts all deformations or all complex threefolds
with the same singular-fibre types.
